\documentclass[12pt,a4paper]{article}

\usepackage[T1]{fontenc}
\usepackage[utf8]{inputenc}
\usepackage[brazil]{babel}
\usepackage{lmodern}
\usepackage{geometry}
\usepackage{hyperref}
\usepackage{enumitem}
\usepackage{array}

\geometry{margin=2.5cm}
\hypersetup{
  colorlinks=true,
  linkcolor=blue,
  urlcolor=blue
}

\title{Plano de Aula: Linguagem SQL --- DDL e DML}
\date{}

\begin{document}
\maketitle

\noindent
\textbf{Duração:} 2 horas (120 minutos)\\
\textbf{Público-alvo:} Iniciantes\\
\textbf{Recursos Necessários:} Laboratório com computadores, SGBD (ex.: PostgreSQL/MySQL) e cliente SQL, projetor (opcional), quadro e marcadores.

\section{Objetivos da Aula}
\begin{itemize}[leftmargin=*,itemsep=0.25em]
  \item Utilizar SQL para implementação de bancos de dados e manipulação de dados.
  \item Aplicar comandos DDL (CREATE/ALTER/DROP) para definição de tabelas.
  \item Aplicar comandos DML (INSERT/UPDATE/DELETE) para manipulação de registros.
  \item Definir PK/FK e restrições básicas na criação das tabelas.
\end{itemize}

\section{Metodologia}
Prática guiada em laboratório. O professor demonstra comandos no projetor e os alunos replicam, com checkpoints curtos. A atividade final envolve criar um pequeno esquema e popular com dados, reforçando integridade referencial.

\section{Cronograma e Conteúdo Detalhado (120 Minutos)}
\subsection*{Parte 1: Preparação do Ambiente (15 minutos)}
\begin{itemize}[leftmargin=*,itemsep=0.35em]
  \item Confirmar acesso ao SGBD e abrir o cliente SQL.
  \item Explicar: banco, schema (quando houver), tabelas, scripts.
\end{itemize}

\subsection*{Parte 2: DDL --- Criando Estrutura (35 minutos)}
\begin{itemize}[leftmargin=*,itemsep=0.35em]
  \item \textbf{CREATE TABLE (20 min):}
  \begin{itemize}[leftmargin=*,itemsep=0.25em]
    \item Criar CLIENTE e PET com PK e FK.
    \item Discutir tipos básicos (texto, número, data) e NOT NULL.
  \end{itemize}
  \item \textbf{ALTER e DROP (15 min):}
  \begin{itemize}[leftmargin=*,itemsep=0.25em]
    \item Alterar coluna (adicionar/renomear) e cuidados com dados existentes.
    \item Quando usar DROP e por que é perigoso.
  \end{itemize}
\end{itemize}

\subsection*{Parte 3: DML --- Inserindo e Alterando Dados (35 minutos)}
\begin{itemize}[leftmargin=*,itemsep=0.35em]
  \item \textbf{INSERT (15 min):} inserir clientes e pets, respeitando FK.
  \item \textbf{UPDATE (10 min):} atualizar telefone/cidade; boas práticas de WHERE.
  \item \textbf{DELETE (10 min):} exclusão e impacto em tabelas relacionadas.
\end{itemize}

\subsection*{Parte 4: Atividade em Duplas --- ``Banco da Clínica'' (25 minutos)}
\begin{itemize}[leftmargin=*,itemsep=0.35em]
  \item Criar um script com:
  \begin{itemize}[leftmargin=*,itemsep=0.25em]
    \item 2 tabelas com PK/FK.
    \item Pelo menos 8 INSERTs.
    \item 2 UPDATEs e 2 DELETEs.
  \end{itemize}
  \item Testar: tentar inserir um PET com ID\_Cliente inexistente e observar o erro.
\end{itemize}

\subsection*{Parte 5: Encerramento (10 minutos)}
\begin{itemize}[leftmargin=*,itemsep=0.35em]
  \item Recapitular: DDL define estrutura; DML manipula dados.
  \item Ponte: ``Próxima aula: SELECTs, filtros, agregações e JOINs.''
\end{itemize}

\end{document}

